\section{Alternatives to ZO0K}

\subsubsection{Proposal 1: VEIL}

In its most simple variant, the followings steps need to be taken.

\subsubsection{Securing the witness queries.}
We need to randomize the witnesses $w_1$ and $w_2$ for the number of queries from $D$. (No R1CS modification necessary). 
Recall we commit the witnesses in interleaved manner, using the FFT decomposition, and thus we need to randomize the witnesses properly, so that each component is affected equally:
Thinking of $\vec w_1$ as the not yet decomposed vector, add a \textit{contiguous} interval of $e \cdot \text{num}_\text{queries}$  random values, where $e$ is the interleaving depth, at any place in the not used padding space of it. 
For $w_2$ likewise, but we may take into account that less than $e$ many components are non-zero.


\subsubsection{Isolator polynomials.}

In its most simple variant, the prover samples addionally to two random polynomials over the extension field $F$,
\[
H_1 \sample F[X_1, \ldots, X_{m_1}], \quad H_2 \sample F[X_1,\ldots, X_{m_2}],
\]
and commits them together with $w_1$ and $w_2$, respectively. 
These `isolator polynomials' are combined together with $w_1$ and $w_2$ in an additional (virtual) batching step before entering WHIR:
\begin{enumerate}
    \item 
    After we combined \eqref{e:phase3:claim:vM} and \eqref{e:phase3:claim:inputs}  into a single inner product claim $\langle K, w\rangle_{H_{m_1 + 1}}$ on the combined witness poly, the prover reveals 
    \[
    v = \langle K, H \rangle_{H_{m_1 + 1}}
    \]
    on the combined isolator polynomial $H = (1 - X_0 ) \cdot H_1 + X_0 \cdot H_2^*$.
    with $w_1$ and $w_2^*$ replaced by $H_1$ and $H_2^*$.  
    \item 
    The verifier draws a further randomness $\gamma_2$ to combine the two claims into a claim on the linear combination $w + \gamma \cdot H$. 
    This combination is not committed, the values of it are computed from those of $w$ and $H$. 
    \item 
    Run WHIR on the virtual linear combination $w + \gamma \cdot H$.
\end{enumerate}

The isolator polynomial entirely decouples the distribution of the future run of WHIR from the past (and thus the witnesses). 
WHIR is run entirely in non-zk mode, no changes here.

While the prove size increases only marginally (by one extension field element per $w_1$ query, likewise for $w_2$), prover cost does, since the prover has to hash the isolator polys. 
Moreover, we most likely hit memory bounds in this variant of VEIL.

\ulrich{%
The witness polys $w_1$ is about $1/4$-th the size as in the BN254 case, and $w_2$ is about $3/4$-th of the size.
However, $H_1$ and $H_2$ together are $3/4$-th of the entire BN254 witnesses, which is why I believe we will hit memory bounds for the largest circuits we want to prove.
}


\paragraph{A more prover friendly variant.}

To mitigate the memory bounds, we insert isolator polynomials in the first round of WHIR instead.
\ulrich{This is a hacky solution, which requires additional sumcheck blinders, I am not a fan of it.}

\begin{enumerate}
\item
The prover only samples a single
\[
H \sample C_1
\]
from the base code of $w_1$ and $w_2^*$, and commits either together with $w_1$. 
This is much less overhead (both compute and memory). 

\item 
Modified first round of WHIR:
\begin{enumerate}
\item
The sumcheck phase is randomized in the usual manner, using 
\[
h_1(X),\ldots, h_{k_1}(X)\sample F[X]^{\leq 3}.
\]
These are committed separately, and in zk manner. 

\item
After specializing the sumcheck claim to $X_1 = \beta_1,\ldots, X_k =\beta_{k_1}$, the prover reveals 
\[
v = \langle K_0(\beta_1,\ldots, \beta_{k_1}, \:.\:), H\rangle_{H_{m_1 - k_1}},
\]
and the claims for $w(\beta_1,\ldots, \beta_{k_1},\:.\:)$ and $H$ are combined via another verifier randomness $\gamma_2\sample F$.
\item 
We now proceed with the virtual linear combination $w(\beta_1,\ldots, \beta_{k_1}, \:.\:) + \gamma_2\cdot H$.
(The query phase, the aggregation phase.)
\end{enumerate}
\end{enumerate}

The value $h_1(\beta_1) +\ldots + h_{k_1}(\beta_{k_1})$ is proven separately. 